\documentclass[10pt]{article}

\usepackage[margin=0.38in]{geometry}
\usepackage{titlesec}
\usepackage{enumitem}
\usepackage{hyperref}
\usepackage[T1]{fontenc}
\usepackage[scaled]{helvet}
\usepackage[none]{hyphenat}
\sloppy

\hypersetup{hidelinks}
\renewcommand{\familydefault}{\sfdefault}
\setlength{\parindent}{0pt}
\setlength{\parskip}{0pt}
\pagestyle{empty}

% Section formatting (ATS-safe, single-column, no graphics)
\setlength{\parskip}{0pt}
\newcommand{\ressection}[1]{
  \vspace{0.25em}
  {\large\bfseries #1}\par
  \vspace{0.08em}\hrule\vspace{0.20em}
}

% Compact bullet layout for one-page engineering resumes
\setlist[itemize]{
  leftmargin=1.2em,
  itemsep=0pt,
  topsep=0pt,
  parsep=0pt,
  partopsep=0pt,
  noitemsep
}

\begin{document}

\begin{center}
{\LARGE {\bfseries Santhosh Sankar}} \\
Ann Arbor, MI \enspace $\vert$ \enspace (734) 846-5255 \enspace $\vert$ \enspace
\href{mailto:ssankar@umich.edu}{ssankar@umich.edu} \enspace $\vert$ \enspace
\href{https://www.linkedin.com/in/santhosh-sankar-3241521b8/}{linkedin.com/in/Tosh-Sankar} \enspace $\vert$ \enspace
\href{https://github.com/ssankar7775}{github.com/ssankar7775}
\end{center}

\ressection{EDUCATION}

{\bfseries University of Michigan}, Ann Arbor, MI \hfill Jan 2026 -- May 2027 (Expected) \\
Master of Engineering in Space Systems \hfill {\bfseries GPA: 4.0/4.0} \\
Climate and Space Sciences and Engineering Department (CLaSP)

\begin{itemize}
  \item Specialization in Controls and Embedded Systems
  \item Relevant Graduate Coursework: Management of Space Systems, Control System Analysis and Design, Astrodynamics
\end{itemize}

{\bfseries University of Michigan}, Ann Arbor, MI \hfill Aug 2021 -- Aug 2024 \\
Bachelor of Science and Engineering in Aerospace Engineering \\
Department of Aerospace Engineering


\ressection{WORK EXPERIENCE}

{\bfseries SpaceX}, Graduate Engineer -- Payload Integration, Launch Engineering \hfill Aug 2025 -- Dec 2025 \\
Vandenberg, CA
\begin{itemize}
\item Engineered an additional payload avionics checkout rack, increasing payload processing throughput by 50\% to support increased launch cadence.
\item Developed a simulation environment to validate automated payload avionics checkout sequences, enabling hardware-independent testing and engineer training.
\item Automated payload checkout test plans, reducing required engineering support by 1 hour per launch campaign.
\item Integrated payload avionics hardware supporting Falcon 9 launch processing operations.
\end{itemize}

\ressection{PROJECT EXPERIENCE}

\textbf{Michigan Experimental Rocketry Organization (MERO)} \hfill Ann Arbor, MI \\
\textit{Avionics Systems Engineer / Responsible Engineer} \hfill Jan 2025 – Present

\begin{itemize}
\item Architecting integrated rocket avionics for flight computer, power distribution, and sensor/actuator interfaces with subsystem-level requirements traceability.
\item Designing custom PCBs for flight computer, power, and I/O breakout hardware supporting guidance, telemetry, and recovery functions.
\item Developing embedded firmware for deterministic sensor acquisition, actuator control, and real-time telemetry handling.
\item Designing and validating vehicle wiring harnesses and end-to-end electrical integration for flight hardware.
\item Architecting LoRa-based flight--ground communication software for command uplink, telemetry downlink, and ground-station monitoring.
\end{itemize}

{\bfseries CubeSat Flight Lab -- University of Michigan}, Systems Engineer \hfill Aug 2023 -- Aug 2024 \\
Ann Arbor, MI
\begin{itemize}
  \item Spearheaded a high-altitude balloon (HAB) flight campaign for a research CubeSat; authored CONOPS, requirements, and integration/test procedures.
  \item Built {\bfseries 6} FlatSat/StratoSat iterations, reducing architecture from 3U 7 kg to 1U 1.2 kg and maturing subsystem integration.
  \item Designed an RP2040-based flight termination unit, reducing subsystem mass by 83\%; validated operation in thermal chamber testing to {\bfseries 110,000 ft altitude}.
  \item Engineered embedded avionics integrating LoRa and GPS; mitigated GPS noise via ground-plane integration for robust command, tracking, and mission safety.
\end{itemize}

{\bfseries Michigan Aeronautical Science Association (MASA)}, Aerodynamics and Recovery Sub-Team Lead \hfill Apr 2022 -- Apr 2023 \\
Ann Arbor, MI
\begin{itemize}
  \item Led aerodynamic and recovery system development for 23 ft. collegiate liquid bi-propellant (RP-1/LOx) rocket Clementine.
  \item Implemented MBSE workflows for the Lonely Mission demonstrator, introducing Rolleron-based R\&D to mitigate roll coupling and improve verification flow.
  \item Engineered a transition from pyrotechnic to pneumatic deployment architecture, lowering fracture risk in fiberglass airframes and improving system reliability.
  \item Optimized fin and fin-can geometry using MATLAB (Mastran), FEA, OpenRocket, and CFD; led manufacturing execution including fabrication and composite layups.
\end{itemize}


\ressection{TECHNICAL SKILLS}

\textbf{Programming:} Python, C/C++, MATLAB, Simulink \\
\textbf{Embedded Systems:} Microcontrollers (RP2040), Raspberry Pi, embedded firmware, sensor interfaces, telemetry \\
\textbf{Avionics \& Instrumentation:} Oscilloscopes, multimeters, DC supplies, thermal chambers, wind tunnels, spectrum analyzers, Logic Analyzers \\
\textbf{Simulation \& Engineering Tools:} Siemens NX, SolidWorks, Star-CCM+, MATLAB, CFD, FEA, Altium Designer \\
\textbf{Hardware Development:} PCB design, wiring harnesses, composite fabrication, aerospace prototyping \\
\textbf{Project Management:} Requirements, test planning, technical documentation, team leadership
\end{document}